\rok{Новембар 2024.}{F}

Први практични тест из Алгоритама и структура података

\zadatak{1}{Quick Sort}
Написати \textit{quick} методу \textit{Quick Sort} алгоритма за сортирање целобројног низа дужине $n$.

\textbf{Додатне напомене за решавање задатка:}
\begin{itemize}
    \item Улаз је несортиран низ целих бројева.
    \item Излаз је сортирани низ целих бројева.
    \item У шаблону се налази класа \texttt{RandomNumberGen} коју треба искористити за генерисање низа случајних бројева.
    \item У оквиру класе \texttt{QuickSort} написати имплементацију за методу:
    \begin{Verbatim}[fontsize=\footnotesize, numbers=left, xleftmargin=10mm]
private static void quickSort(int[] arr, int lower, int upper)
    \end{Verbatim}
    \item Обезбедити код у \texttt{Main} методи за тестирање алгоритма.
\end{itemize}



\zadatak{2}{Брисање непарних бројева из стека, производ преосталих}
Имплементирати методу која прима низ целих бројева. Метода треба да користи структуру података стека (класа \texttt{Stack}) за организацију и манипулацију елементима, и треба да изврши следеће кораке:
\begin{itemize}
    \item Уклањање непарних бројева: пролазом кроз стек, уклонити све бројеве који су непарни.
    \item Рачунање производа преосталих бројева: израчунати производ свих бројева који су остали у стеку након уклањања непарних бројева.
    \item Враћање резултата: као резултат, метода треба да врати израчунати производ преосталих бројева у стеку.
\end{itemize}

\textbf{Додатни захтеви за решавање задатка:}
\begin{itemize}
    \item Задатак решавати помоћу структуре стека (класа \texttt{Stack}) која је дата у оквиру шаблона задатка.
    \item Задатак решавати у оквиру методе \texttt{public int RemoveOddAndReturnProduct (int[] array)}.
    \item Након што се уклоне непарни бројеви из стека и израчуна производ преосталих бројева, стек треба да остане у истом стању као пре почетног процеса. То значи да преостали бројеви морају бити враћени у стек у оригиналном редоследу, а сви измењени елементи (они који су уклоњени) не смеју се вратити.
    \item Алгоритам је неопходно написати са линеарном временском комплексношћу $O(n)$.
\end{itemize}

\textbf{Потпис методе:}
\begin{Verbatim}[fontsize=\small]
public int RemoveOddAndReturnProduct (int[] array)
\end{Verbatim}

\textbf{Примери:}
\begin{itemize}
    \item \textbf{Улаз:} \texttt{[10, 15, 7, 30, 2, 8]} \\
          \textbf{Излаз:} Производ преосталих бројева: \texttt{4800}; преостали бројеви у стеку: \texttt{10, 30, 2, 8}
    \item \textbf{Улаз:} \texttt{[5, 2, 8, 3, 4, 6]} \\
          \textbf{Излаз:} Производ преосталих бројева: \texttt{384}; преостали бројеви у стеку: \texttt{2, 8, 4, 6}
\end{itemize}



\zadatak{3}{Разлика максималне вредности локалних максимума и минималне вредности локалних минимума}
Имплементирати методу која прима повезану листу целих бројева и проналази све локалне максимуме и локалне минимуме, затим израчунава разлику између максималне вредности локалних максимума и минималне вредности локалних минимума.

\textbf{За дефиницију:}
\begin{itemize}
    \item Локални максимум је број који је већи од претходног и следећег броја.
    \item Локални минимум је број који је мањи од претходног и следећег броја.
\end{itemize}

Ако нема локалних максимума или минимума, метода треба да врати -1.

\textbf{Додатни захтеви за решавање задатка:}
\begin{itemize}
    \item Задатак решавати помоћу структуре листе (класа \texttt{LinkedList}) која је дата у оквиру шаблона задатка.
    \item Не користити додатне колекције попут листе или низа, користити само повезану листу.
    \item Захтеван потпис методе:
    \begin{Verbatim}[fontsize=\footnotesize, numbers=left, xleftmargin=10mm]
public int findLocalMaxMinDifference()
    \end{Verbatim}
\end{itemize}

\textbf{Пример 1:}
\begin{itemize}
    \item \textbf{Улаз:} \texttt{[1, 3, 2, 5, 6, 4, 7, 3]}
    \item Локални максимуми: \texttt{3, 6, 7} (индекси 1, 4, 6)
    \item Локални минимуми: \texttt{2, 4} (индекси 2, 5)
    \item Максимална вредност локалних максимума је \texttt{7}, а минимална вредност локалних минимума је \texttt{2}.
    \item \textbf{Излаз:} \texttt{5}
\end{itemize}

\noindent\textbf{Пример 2:}
\begin{itemize}
    \item \textbf{Улаз:} \texttt{[1, 2, 3, 4, 5]}
    \item Нема локалних минимума и максимума.
    \item \textbf{Излаз:} \texttt{-1}
\end{itemize}

\sablon{AISP2024_Test1_GrupaF}{2024/Novembar/GrupaF/AISP2024_Test1_GrupaF.linq}
